% FIG ch06-02 flow grammar (Chinese)%
% Shared styles for chapter 6 workflow TikZ figures.%
% Every figure inputs this file, so commands are defined with \providecommand.%
\providecolor{BookTeal}{HTML}{176C68}%
\providecolor{BookInk}{HTML}{23363E}%
\providecolor{BookRule}{HTML}{C8D6D6}%
\providecolor{BookPale}{HTML}{F1F6F5}%
\definecolor{WetPale}{HTML}{FBF1EA}%
\definecolor{Wet}{HTML}{B5653A}%
\definecolor{Seq}{HTML}{8A6D1F}%
\definecolor{DryPale}{HTML}{EEF4F7}%
\definecolor{Dry}{HTML}{2F6F8F}%
\definecolor{ChSixFeedback}{HTML}{7A8A90}%
\definecolor{ChSixBad}{HTML}{C0392B}%
\usetikzlibrary{arrows.meta,positioning,shapes.geometric,shapes.misc,fit,backgrounds,calc}%
\ifdefined\ChSixLayersReady\else%
  \pgfdeclarelayer{chsixlane}%
  \pgfsetlayers{chsixlane,background,main}%
  \def\ChSixLayersReady{}%
\fi%

\tikzset{%
  chsix/.style={font=\sffamily\scriptsize, text=BookInk, align=center},%
  flowstep/.style={%
    chsix, rounded corners=2pt, draw=Wet, line width=0.45pt, fill=white,%
    text width=19mm, minimum height=7.4mm, inner sep=1.3pt,%
  },%
  flowstepdry/.style={flowstep, draw=Dry},%
  flowstepseq/.style={flowstep, draw=Seq},%
  artifact/.style={%
    chsix, trapezium, trapezium left angle=70, trapezium right angle=110,%
    trapezium stretches=true,%
    draw=BookInk, line width=0.45pt, fill=white,%
    minimum height=6.2mm, inner sep=1.2pt,%
  },%
  gate/.style={%
    chsix, diamond, aspect=1.7, draw=BookInk, line width=0.45pt, fill=white,%
    text width=12mm, inner sep=0.6pt, minimum height=8mm,%
  },%
  opt/.style={flowstep, dashed, draw=BookInk!70},%
  optdry/.style={flowstepdry, dashed, draw=BookInk!70},%
  stack/.style={flowstep},%
  link/.style={-{Stealth[length=1.6mm]}, semithick, draw=BookInk},%
  linkthin/.style={-{Stealth[length=1.2mm]}, thin, draw=BookInk!80},%
  back/.style={-{Stealth[length=1.5mm]}, dashed, draw=ChSixFeedback},%
  conn/.style={%
    circle, draw=BookInk, line width=0.5pt, minimum size=4.6mm,%
    font=\sffamily\scriptsize\bfseries, inner sep=0pt, text=BookInk,%
  },%
  lanetitle/.style={font=\sffamily\bfseries\small, text=BookInk, align=left},%
  passlabel/.style={font=\sffamily\scriptsize, text=BookInk, fill=white, inner sep=0.7pt, align=center},%
  mergelabel/.style={font=\sffamily\scriptsize, text=BookInk, fill=white, inner sep=0.8pt, align=center},%
  taglab/.style={font=\sffamily\scriptsize, text=ChSixFeedback, fill=white, inner sep=0.8pt, align=center},%
  gramcap/.style={font=\sffamily\scriptsize, text=BookInk, align=left, anchor=west},%
}%

% Offset card behind a two-sample step. Text stays on the front node only.%
\providecommand{\ChSixCardBack}[1]{%
  \begin{pgfonlayer}{background}%
    \draw[Wet, line width=0.45pt, rounded corners=2pt, fill=white]%
      ([shift={(-0.55mm,0.55mm)}]#1.south west)%
      rectangle ([shift={(-0.55mm,0.55mm)}]#1.north east);%
  \end{pgfonlayer}%
}%
% #1 options, #2 name, #3 position (with parentheses), #4 text%
\providecommand{\ChSixCard}[4]{%
  \node[stack, #1] (#2) at #3 {#4};%
  \ChSixCardBack{#2}%
}%

\providecommand{\ChSixMaybeLabel}[1]{%
  \ifdefined\nolabels\else#1\fi%
}%

\providecommand{\ChSixIcon}[2]{%
  \ifdefined\ChSixShowIcons%
    \IfFileExists{#1}{%
      \node[anchor=south, inner sep=0] at ([yshift=0.7mm]#2.north)%
        {\includegraphics[height=5.4mm]{#1}};%
    }{}%
  \fi%
}%
%
\begin{tikzpicture}[x=1mm, y=1mm, font=\sffamily\scriptsize, text=BookInk]%
  \node[anchor=north west, font=\sffamily\bfseries\small] (titleL) at (0,0) {节点与箭头语法};%
  \draw[BookRule, line width=0.4pt] (0,-7) rectangle (128,-126);%

  \node[flowstep] (g1) at (16,-14) {操作步骤};%
  \node[gramcap] at (36,-14) {圆角矩形——做什么};%

  \node[artifact] (g2) at (16,-25) {产物};%
  \node[gramcap] at (36,-25) {平行四边形——得到什么、交给下一步什么};%

  \node[gate] (g3) at (16,-38) {质控};%
  \node[passlabel, anchor=west] at ([xshift=1.5mm]g3.east) {通过};%
  \node[taglab, anchor=north west] at ([xshift=1.5mm, yshift=-0.6mm]g3.south east) {不通过};%
  \node[gramcap] at (36,-38) {菱形——须有两个出口（通过 / 不通过）};%

  \node[opt] (g4) at (16,-50) {可选步骤};%
  \node[gramcap] at (36,-50) {虚线圆角矩形};%

  \node[stack] (gs1) at (13,-61) {};%
  \ChSixCardBack{gs1}%
  \node[stack] (gs2) at (16.2,-61) {每样本};%
  \node[gramcap] at (36,-61) {错位叠放——每个样本分别进行};%

  \draw[link] (6,-72) -- ++(16,0);%
  \node[gramcap] at (36,-72) {实线箭头——顺序与交接};%

  \draw[back] (6,-83) to[out=20, in=160] (24,-83);%
  \node[gramcap] at (36,-83) {虚线弧——反馈或返回};%

  \fill (16,-94) circle (0.7pt);%
  \node[gramcap] at (36,-94) {实心圆点——分叉 / 合并点};%

  \node[conn] (gc) at (16,-105) {A};%
  \node[gramcap] at (36,-105) {圆形连接符——跨行续接};%

  \fill[WetPale, draw=Wet, line width=0.3pt] (6,-116) rectangle (16,-122);%
  \fill[DryPale, draw=Dry, line width=0.3pt] (18,-116) rectangle (28,-122);%
  \node[gramcap] at (36,-119) {泳道底色——湿实验 / 计算};%

  \node[anchor=north west, font=\sffamily\bfseries\small] (titleR) at (0,-130) {流程箭头 ≠ 因果};%
  \node[flowstep] (d1) at (18,-144) {组织解离};%
  \node[artifact] (d2) at (46,-144) {单细胞悬液};%
  \draw[link] (d1) -- (d2);%
  \node[gramcap] at (68,-144) {流程：先后与交接};%

  \node[flowstep] (e1) at (18,-160) {组织解离};%
  \node[gramcap] at (40,-160) {$\Rightarrow$ 应激相关基因表达升高};%
  \draw[BookInk!45, densely dotted] (e1.east) -- ++(6,0);%
  \node[gramcap, anchor=north west, text width=120mm] at (0,-170)%
    {机制或因果信息写成注释或图注，不画成流程箭头（van den Brink 等 2017）。};%
\end{tikzpicture}%
